\cotefile{128}
\scanpage{17}
2°) $A'$ \uline{est radiciel sur} $A$, i.e. $\operatorname{Ker}(A'\otimes_{A} A'\to A')$
est nilpotent, donc $\Delta$ \uline{est nilpotent}. Alors
posons
\[
\Omega=\Delta/\Delta^2
\]
et distinguons deux cas

a) $\Omega=0$ d'où $\Delta=\Delta^2$, \deleted{puisque} et $\Delta=0$,
puisque $\Delta$ est nilpotent, donc $A''\simeq A'$, donc
$A^{\unsure{\Delta|1}}\simeq\operatorname{Ker}(A'\rightrightarrows A'')=A'$, donc $\operatorname{long} A'/A=0$.

b) $\Omega\ne0$ ; Notons \unsure{qu'en|que dans} tous cas, \unsure{posons|posant}
pour $x'\in A'$
\[
\begin{aligned}
d(x')&=\delta(x')\bmod\Delta^2\quad\in\Omega\\
&=\varphi(d_{A'/A}(x'))
\end{aligned}
\]
\unsure{l'app.|lorsq.|bien}
où $\varphi:\Omega^1_{A'/A}\to\Omega$ est l'application canonique,
$d$ est une $A$-dérivation de $A'$ dans $\Omega$ :
\[
d(x'y')=x'dy'+y'dx'
\]
d'où
\[
d(\mathfrak m'{}^2)\subset\mathfrak m'\Omega
\]
et comme $\Omega\ne0$, \deleted{\unsure{on} $\mathfrak m'\Omega\ne\Omega$ par Naka} \inserted{un élément de $\mathfrak m'\Omega$ ne}
peut être un générateur de $\Omega$ par Nakayama,
d'où d'après \deleted{$\mathfrak m$} l'hypothèse sur $A'/A$.
\[
\mathfrak m'{}^2\subset A\qquad\text{d'où }\mathfrak m'{}^2\subset\mathfrak m \tag*{$(*)$}
\]
Or comme $A'$ est local à corps résiduels
% archive sheet number: 16
